\oldpage[IV-3]{173}
\section{Étude des fibres des morphismes plats de présentation finie}
\label{section:IV.12-fr}

\setcounter{subsection}{-1}
\subsection{Introduction}
\label{subsection:IV.12.0-fr}

\emph{Nous utiliserons dans tout ce paragraphe les notations générales de}
(9.4.1).

\begin{env}[12.0.1]
\label{IV.12.0.1-fr}
Étant donné un morphisme localement de présentation finie $f:X\to Y$, nous
avons vu (9.9) que pour certaines propriétés locales $\mathbf P$ de préschémas
sur des corps ou de Modules sur ces préschémas, l'ensemble $E$ des $x\in X$
tels que la propriété $\mathbf P$ ait lieu, pour la fibre $X_{f(x)}$, au point
$x$ de cette fibre, est localement constructible dans $X$. Nous nous proposons
de montrer que, pour la plupart de ces propriétés, si l'on suppose de plus que
le morphisme $f$ est plat, alors l'ensemble $E$ est même ouvert dans $X$. De
même ((9.2) à (9.8)) nous avons montré que si $f$ est de présentation finie,
et si cette fois $\mathbf P$ désigne certaines propriétés globales de
préschémas sur des corps ou de Modules sur ces préschémas, l'ensemble $F$ des
$y\in Y$ tels que la propriété $\mathbf P$ ait lieu pour la fibre $X_y$ est
localement constructible dans $Y$. Nous montrerons que si l'on suppose de plus
que le morphisme $f$ est propre et plat, $F$ est même ouvert dans $Y$.
\end{env}

\begin{env}[12.0.2]
\label{IV.12.0.2-fr}
La méthode générale de démonstration des propriétés en question comporte trois
étapes. On se ramène d'abord au cas où $Y$ est affine et $X$ de présentation
finie~; puis~:
\begin{enumerate}[label=\emph{\Alph*)}]
  \item À l'aide de (8.9.1) (et éventuellement d'autres résultats du \S~8) et
  de (11.2.6), on se ramène au cas où $X$ et $Y$ sont noethériens.
  \item On applique les résultats du \S~9 rappelés en (12.0.1) prouvant que
  $E$ (resp. $F$) est constructible.
  \item Pour voir que $E$ est ouvert, il suffit, en vertu de
  ($\mathbf 0_{\mathrm{III}}$, 9.2.5) de montrer que si $x\in E$, alors toute
  générisation $x'$ de $x$ appartient aussi à $E$. Utilisant ($\mathbf{II}$,
  7.1.7), on voit, puisque $Y$ est noethérien, qu'il existe un spectre d'anneau
  de valuation discrète $Y_1$ et un morphisme $h:Y_1\to X$ tel que, si $y_1$
  (resp. $y'_1$) est le point fermé (resp. le point générique) de $Y_1$, on ait
  $h(y_1)=x$, $h(y'_1)=x'$. On fait alors le changement de base
  $g=f\circ h:Y_1\to Y$~; compte tenu des résultats des \S\S~4 et 6 sur les
  préschémas localement noethériens sur des corps et les changements du corps
  de base, on est ramené à prouver l'assertion en question pour un point $x_1$
  de $X_1=X\times_YY_1$ au-dessus de $y_1$ et pour une générisation $x'_1$ de
  $x_1$ au-dessus de $y'_1$. Comme $Y_1=\operatorname{Spec}(A)$, où $A$ est
  un anneau de valuation discrète, d'uniformisante $t$, on doit en définitive,
  pour un $A$-module $M$, prouver une propriété de $M$, sachant que $M/tM$ a
  la même propriété et que $t$ est $M$-régulier (ce qui résulte de l'hypothèse
  de platitude)~; on utilise pour cela les résultats de (3.4) et de (5.12). On
  procède de même pour l'ensemble $F\subset Y$.
\end{enumerate}
\end{env}

\oldpage[IV-3]{174}
\subsection{Propriétés locales des fibres d'un morphisme plat localement de
présentation finie}
\label{subsection:IV.12.1-fr}

\begin{theorem}[12.1.1]
\label{IV.12.1.1-fr}
Soient $f:X\to Y$ un morphisme localement de présentation finie, $\mathcal F$
un $\mathcal O_X$-Module $f$-plat et de présentation finie, $\Phi$ une partie
finie de $\mathbf Z\cup\{\pm\infty\}$, $k$ un entier. Les parties suivantes de
$X$ sont ouvertes~:
\begin{enumerate}[label=(\roman*)]
  \item L'ensemble des $x\in X$ tels que les dimensions des cycles premiers
  associés à $\mathcal F_{f(x)}$ et contenant $x$ soient des éléments de
  $\Phi$.
  \item L'ensemble des $x\in X$ tels que les cycles premiers associés à
  $\mathcal F_{f(x)}$ contenant $x$ aient tous la dimension $k$.
  \item L'ensemble des $x\in X$ tels que $x$ n'appartienne à aucun cycle
  premier immergé associé à $\mathcal F_{f(x)}$.
  \item L'ensemble des $x\in X$ tels que $\mathcal F_{f(x)}$ soit
  équidimensionnel au point $x$ et possède la propriété $(S_k)$ au point $x$
  (5.7.2).
  \item L'ensemble des $x\in X$ tels que
  $\operatorname{coprof}((\mathcal F_{f(x)})_x)\leq k$
  ($\mathbf 0$, 16.4.9).
  \item L'ensemble des $x\in X$ tels que $\mathcal F_{f(x)}$ soit un
  $\mathcal O_{X_{f(x)}}$-Module de Cohen-Macaulay au point $x$ (5.7.1).
  \item L'ensemble des $x\in X$ tels que $\mathcal F_{f(x)}$ soit
  géométriquement réduit au point $x$ (4.6.22).
  \item L'ensemble des $x\in X$ tels que $\mathcal F_{f(x)}$ soit
  géométriquement ponctuellement intègre au point $x$ (4.6.22).
\end{enumerate}
\end{theorem}

\begin{proof}
Les questions étant locales sur $X$, on se ramène d'abord au cas où
$X=\operatorname{Spec}(B)$ et $Y=\operatorname{Spec}(A)$ sont affines, $f$ un
morphisme de présentation finie. Il y a alors un sous-anneau noethérien $A_0$
de $A$, un $A_0$-préschéma de type fini $X_0$, un
$\mathcal O_{X_0}$-Module cohérent $\mathcal F_0$ tel que
$\mathcal F_0\otimes_{A_0}A$ soit isomorphe à $\mathcal F$ (8.9.1)~; en
outre, en vertu de (11.2.6), on peut supposer que $\mathcal F_0$ est
$Y_0$-plat (avec $Y_0=\operatorname{Spec}(A_0)$). Si $h:X\to X_0$ est la
projection canonique, l'ensemble $E$ des points $x\in X$ où l'une des
propriétés (i) à (viii) est vérifiée est égal à $h^{-1}(E_0)$, où $E_0$ est
l'ensemble des $x_0\in X_0$ où la propriété correspondante pour $\mathcal F_0$
est vérifiée~: cela résulte, pour les propriétés (i) à (iii), de (4.2.7)~; pour
les propriétés (iv) à (vi), de (6.7.1)~; pour les propriétés (vii) et (viii),
de (4.7.11).

On est ainsi ramené au cas où $A$ et $B$ sont noethériens, $f$ de type fini,
et $\mathcal F=\widetilde M$, où $M$ est un $B$-module de type fini. En vertu
de (9.9.2), on sait que l'ensemble $E$ est constructible pour toutes les
propriétés envisagées, et il reste dans chaque cas l'étape C) de (12.0.2), où
l'on doit prouver que $E$ est stable par générisation.

\begin{env}[12.1.1.1]
\label{IV.12.1.1.1-fr}
Commençons par le cas (iii) qui est le plus simple. Soit $x'\ne x$ une
générisation de $x$ dans $X$. Posons $y=f(x)$, $y'=f(x')$~; il existe un
spectre d'anneau de valuation discrète $Y_1$ et un morphisme $u:Y_1\to X$ tels
que, si $y_1$ et $y'_1$ sont le point fermé et le point générique de $Y_1$, on
ait $u(y_1)=x$ et $u(y'_1)=x'$ ($\mathbf{II}$, 7.1.7)~; posons
$g=f\circ u:Y_1\to Y$~; si $X_1=X\times_YY_1$, il existe une $Y_1$-section
$u_1:Y_1\to X_1$ telle que $u=g_1\circ u_1$, où $g_1:X_1\to X$ est la
projection canonique. Si $x_1=u_1(y_1)$, $x'_1=u_1(y'_1)$, on a donc
$g_1(x_1)=x$ et $g_1(x'_1)=x'$, et $x'_1$ est une générisation de $x_1$.
Appliquant de nouveau (4.2.7), on voit
\oldpage[IV-3]{175}
qu'on est ramené au cas où $Y=\operatorname{Spec}(A)$ est le spectre d'un
anneau de valuation discrète, $y$ étant le point fermé et $y'$ le point
générique de $Y$. Si $t$ est une uniformisante de $A$, l'hypothèse que
$\mathcal F$ est $f$-plat entraîne que $t$ est $\mathcal F$-régulier
($\mathbf 0_{\mathrm I}$, 6.3.4). Par hypothèse aucun des cycles premiers
associés immergés de $\mathcal F/t\mathcal F$ ne contient $x$. Alors, il
résulte de (3.4.4) que $x$ n'appartient à aucun cycle premier immergé associé à
$\mathcal F$, et il en est donc de même de toute générisation de $x$. En
particulier, $x'$ n'appartient à aucun cycle premier immergé associé à
$\mathcal F$, ni \emph{a fortiori} à aucun de ceux associés à
$\mathcal F_{y'}$ ($X_{y'}$ étant un sous-préschéma induit sur un ouvert de
$X$).
\end{env}

\begin{env}[12.1.1.2]
\label{IV.12.1.1.2-fr}
Envisageons ensuite les cas (iv) et (v) ((vi) n'est qu'un cas particulier de
(v)). On procède comme ci-dessus (en utilisant (6.7.1) au lieu de (4.2.7)) et
l'on est ramené au cas où $Y$ est le spectre d'un anneau de valuation discrète
$A$.

L'anneau $C=\mathcal O_{X,x}$ est alors localisé d'une $A$-algèbre de type
fini, donc est caténaire (5.6.4), et par hypothèse le $C$-module
$M_x/tM_x$ vérifie $(S_k)$ et est équidimensionnel (resp. on a
$\operatorname{coprof}(M_x/tM_x)\leq k$)~; on déduit donc de (5.12.2) (resp.
de ($\mathbf 0$, 16.4.10)) que $M_x$ vérifie $(S_k)$ et est équidimensionnel
(resp. que $\operatorname{coprof}(M_x)\leq k$) puisque $t$ est $M_x$-régulier
et appartient à l'idéal maximal de $C$. D'où les conclusions puisque $x'$ est
une générisation de $x$ et que $(\mathcal F_{y'})_{x'}=\mathcal F_{x'}$.
\end{env}

\begin{env}[12.1.1.3]
\label{IV.12.1.1.3-fr}
Passons aux cas (vii) et (viii). On peut évidemment remplacer $Y$ par le
sous-schéma intègre ayant $\overline{\{y'\}}$ pour espace sous-jacent, et $X$
par $f^{-1}(\overline{\{y'\}})$, sans changer les fibres en $y$ et $y'$, donc
on peut supposer $A$ intègre, de corps des fractions $K=k(y')$. On sait
((4.5.11) et (4.7.8)) qu'il existe une extension finie $K'$ de $K$ telle que,
pour le $(\mathcal O_X\otimes_AK')$-Module
$\mathcal F'=\mathcal F\otimes_AK'$, les cycles premiers associés soient
géométriquement irréductibles et que les longueurs géométriques de
$\mathcal F'$ aux points maximaux $z'_i$ de son support soient respectivement
égales aux longueurs de $\mathcal F'_{z'_i}$ en ces points~; il revient donc
au même de dire que $\mathcal F$ est géométriquement réduit (resp.
géométriquement ponctuellement intègre) en un point $x'$ de
$X_{y'}=X\otimes_AK$, ou de dire que $\mathcal F'$ est réduit (resp. intègre)
aux points de $X\otimes_AK'$ au-dessus de $x'$, compte tenu de (4.2.7). Soit
$A'$ une $A$-algèbre finie dont $K'$ est le corps des fractions~; il résulte
de (4.7.11) qu'en tout point de $X\otimes_AA'$ au-dessus de $x$,
$\mathcal F\otimes_AA'$ a la propriété (vii) (resp. (viii)). Remplaçant $A$
par $A'$ et $X$ par $X\otimes_AA'$, on voit qu'on peut se borner à prouver que
$\mathcal F$ est réduit (resp. intègre) en toute générisation $x'$ de $x$. On
procède alors comme dans (12.1.1.1), en utilisant cette fois (4.7.11), et l'on
est encore ramené au cas où $A$ est un anneau de valuation discrète, $y$ étant
le point fermé, $y'$ le point générique de $Y$, et l'uniformisante $t$ de $A$
étant un élément $M$-régulier. Comme par hypothèse
$\mathcal F/t\mathcal F$ est réduit (resp. intègre) au point $x$, il résulte de
(3.4.6) (resp. (3.4.5)) que $\mathcal F$ est réduit (resp. intègre) au point
$x$, donc aussi dans un voisinage de $x$, et en particulier au point $x'$, ce
qui achève la démonstration dans les cas (vii) et (viii).
\end{env}

\begin{env}[12.1.1.4]
\label{IV.12.1.1.4-fr}
Reste à examiner les cas (i) et (ii). Remplaçant $Y$ par le sous-schéma intègre
ayant $\overline{\{y'\}}$ pour espace sous-jacent, on peut se borner au cas où
$Y$ est irréductible et $y'$ son point générique. Soit $z'$ un point générique
d'un cycle premier associé à $\mathcal F_{y'}$
\oldpage[IV-3]{176}
contenant $x'$, et soit $Z'$ l'adhérence de $z'$ dans $X$, de sorte que l'on a
$x\in Z'$. Pour traiter le cas (i), il va suffire de prouver la

\begin{proposition}[12.1.1.5]
\label{IV.12.1.1.5-fr}
Soient $Y$ un préschéma localement noethérien, $f:X\to Y$ un morphisme
localement de type fini, $\mathcal F$ un $\mathcal O_X$-Module cohérent et
$f$-plat~; pour tout $y\in Y$, posons
$\mathcal F_y=\mathcal F\otimes_{\mathcal O_Y}k(y)$. Soient $z'$ un point de
$X$, $y'=f(z')$, et supposons que $z'\in\operatorname{Ass}(\mathcal F_{y'})$.
Soient $Z'$ (resp. $Y'$) le sous-préschéma réduit de $X$ (resp. $Y$) ayant
pour espace sous-jacent l'adhérence de $\{z'\}$ dans $X$ (resp. de $\{y'\}$
dans $Y$). Alors, pour tout $y\in f(Z')$, les dimensions de toutes les
composantes irréductibles de $Z'_y$ sont égales à $\dim(Z'_{y'})$ (dimension
de l'adhérence de $z'$ dans $X_{y'}$) (ce que nous exprimerons plus tard
(13.2.2) en disant que la restriction $Z'\to Y'$ de $f$ est un morphisme
équidimensionnel)~; en outre, en tout point maximal $z$ de $Z'_y$, on a
$z\in\operatorname{Ass}(\mathcal F_y)$.
\end{proposition}

Pour cela, nous allons nous ramener au cas où $Y$ est spectre d'un anneau de
valuation discrète. Prenons un spectre d'anneau de valuation discrète $Y_1$ et
un morphisme $h:Y_1\to X$ tel que $h(y_1)=z$, $h(y'_1)=z'$, $y_1$ et $y'_1$
étant le point fermé et le point générique de $Y_1$ respectivement
($\mathbf{II}$, 7.1.7). Soient $g=f\circ h:Y_1\to Y$,
$X_1=X\times_YY_1$, et soient $f_1:X_1\to Y_1$, $g_1:X_1\to X$ les projections
canoniques~; il y a une $Y_1$-section $h_1:Y_1\to X_1$ telle que
$h=g_1\circ h_1$ ($\mathbf I$, 3.3.14). Si
$Z'_1=Z'_{y'}\otimes_{k(y')}k(y'_1)$, on sait (4.2.6) que les composantes
irréductibles de $Z'_1$ dominent $Z'_{y'}$~; soit $z'_1$ le point générique de
l'une de ces composantes qui contient $h_1(y'_1)$, de sorte que
$f_1(z'_1)=y'_1$ et $g_1(z'_1)=z'$. Comme $h_1(y_1)$ est spécialisation de
$h_1(y'_1)$, il est \emph{a fortiori} spécialisation de $z'_1$~; soit $z_1$ un
point générique d'une des composantes irréductibles de
$\overline{\{z'_1\}}\cap(X_1)_{y_1}$ contenant $h_1(y_1)$. Alors $g_1(z_1)$
est spécialisation de $g_1(z'_1)=z'$ dans $X$, et
$z=h(y_1)=g_1(h_1(y_1))$ est spécialisation de $g_1(z_1)$~; mais comme
$f(g_1(z_1))=y$ et que $z$ est point générique de $Z'_y$, on a nécessairement
$z=g_1(z_1)$.

Supposons maintenant que l'on ait prouvé que, si l'on pose
$\mathcal F_1=\mathcal F\otimes_{\mathcal O_Y}\mathcal O_{Y_1}$, on a
$z_1\in\operatorname{Ass}((\mathcal F_1)_{y_1})$~; il résultera de (4.2.7) que
$z\in\operatorname{Ass}(\mathcal F_y)$~; en outre, les dimensions de
$\overline{\{z_1\}}\cap(X_1)_{y_1}$ et de $\overline{\{z\}}\cap X_y$ sont
égales, et de même les dimensions de
$\overline{\{z'_1\}}\cap(X_1)_{y'_1}$ et de
$\overline{\{z'\}}\cap X_{y'}$ (4.2.7)~; on s'est donc bien ramené, comme
annoncé, au cas où $Y$ est spectre d'un anneau de valuation discrète $A$, $y$
et $y'$ étant son point fermé et son point générique respectivement.

Cela étant, comme $z'\in\operatorname{Ass}(\mathcal F_{y'})$, on a
\emph{a fortiori} $z'\in\operatorname{Ass}(\mathcal F)$. Si $t$ est une
uniformisante de $A$, $t$ est $\mathcal F$-régulier par platitude, donc (3.4.3)
on a $z\in\operatorname{Ass}(\mathcal F/t\mathcal F)
=\operatorname{Ass}(\mathcal F_y)$. Quant à l'assertion relative aux
dimensions, elle résulte de (7.1.13), appliqué à un sous-préschéma fermé de $X$
ayant $Z'$ pour espace sous-jacent.

Abordons enfin le cas (ii) de (12.1.1). Avec les mêmes notations, il faut
prouver que si tous les cycles premiers associés à $\mathcal F_y$ et contenant
$x$ ont la dimension $k$, il en est de même de tous les cycles premiers
associés à $\mathcal F_{y'}$ et contenant $x'$. Or on vient de voir que tout
cycle premier associé à $\mathcal F_{y'}$ et contenant $x'$ a même dimension
que l'un des cycles premiers associés à $\mathcal F_y$ et contenant $x$~; cela
démontre donc (ii).
\end{env}
\end{proof}

\begin{remarks}[12.1.2]
\label{IV.12.1.2-fr}
\begin{enumerate}[label=(\roman*)]
  \item Sous les conditions de (12.1.1.5) avec $\mathcal F=\mathcal O_X$ (de
  sorte que $f$ est plat), on ne peut en général affirmer que la restriction
  $Z'\to Y'$ de $f$ est un morphisme plat ni même un morphisme ouvert. C'est ce
  que montre l'exemple (6.5.5,
  \oldpage[IV-3]{177}
  (ii)) où l'on prend pour $Z'$ une des deux composantes irréductibles de
  $X=\operatorname{Spec}(B)$. Il est immédiat que ce morphisme restriction
  n'est pas ouvert aux points de $Z'$ au-dessus du «~point double~» de $Y'=Y$.
  \item Dans les hypothèses de (12.1.1.5), avec
  $\mathcal F=\mathcal O_X$, on ne peut affaiblir la condition «~$f$ est
  plat~» en «~$f$ est universellement ouvert~» (cf. (2.4.6)), comme nous le
  verrons plus loin sur un exemple (14.4.10, (i)).
\end{enumerate}
\end{remarks}

\begin{corollary}[12.1.3]
\label{IV.12.1.3-fr}
Sous les hypothèses de (12.1.1), la fonction
$x\mapsto\operatorname{coprof}((\mathcal F_{f(x)})_x)$ est semi-continue
supérieurement dans $X$ et la fonction
$x\mapsto\operatorname{prof}_x^*(\mathcal F_{f(x)})$ (10.8.1) est
semi-continue inférieurement dans $X$.
\end{corollary}

\begin{proof}
La première assertion n'est autre qu'une formulation équivalente de (12.1.1,
(v)). Pour prouver la seconde, on peut d'abord, compte tenu de (10.8.7) et
(10.8.8), se ramener comme dans (12.1.1) au cas où $Y$ est le spectre d'un
anneau de valuation discrète $A$ de point fermé $y$ et de point générique $y'$,
$X=\operatorname{Spec}(B)$ étant affine, $\mathcal F=\widetilde M$, où $M$ est
un $B$-module de type fini. Puisque $\mathcal F$ est $f$-plat, toute composante
irréductible de $\operatorname{Supp}(\mathcal F)$ qui contient $x\in X_y$
domine $Y$ (2.3.4), autrement dit son point générique $z'$ appartient à
$X_{y'}$~; les composantes irréductibles $Z$ de
$\operatorname{Supp}(\mathcal F)$ qui contiennent une générisation $x'$ de
$x$ appartenant à $X_{y'}$ sont donc exactement celles qui contiennent $x$ et
en outre, d'après (12.1.1.5), on a $\dim(Z_{y'})=\dim(Z_y)$, d'où, si
$T=\operatorname{Supp}(\mathcal F)$,
$\dim_x(T_y)=\dim_{x'}(T_{y'})$. De plus, pour une uniformisante $t$ de $A$, on
a vu dans (12.1.1, (v)) que l'on a
$\operatorname{coprof}(M_x)=\operatorname{coprof}(M_x/tM_x)$, et puisque
d'autre part $\operatorname{coprof}(M_{x'})\leq\operatorname{coprof}(M_x)$
(6.11.5), on voit que l'on a
$\operatorname{coprof}((\mathcal F_{y'})_{x'})\leq
\operatorname{coprof}((\mathcal F_y)_x)$, la relation
$\operatorname{prof}_x^*(\mathcal F_y)\leq
\operatorname{prof}_{x'}^*(\mathcal F_{y'})$ découle alors de (10.8.7).
\end{proof}

\begin{remark}[12.1.4]
\label{IV.12.1.4-fr}
On déduit aussi de (12.1.1, (i)) que la fonction
$x\mapsto\dim_x(X_{f(x)})$ est semi-continue supérieurement dans $X$, mais
nous verrons plus loin (13.1.3) que cette propriété est vraie même sans
supposer $f$ plat.
\end{remark}

\begin{corollary}[12.1.5]
\label{IV.12.1.5-fr}
Soient $Y$ un préschéma localement noethérien, $f:X\to Y$ un morphisme
localement de type fini, $\mathcal F$ un $\mathcal O_X$-Module cohérent et
$f$-plat, $\mathcal G$ un $\mathcal O_Y$-Module cohérent, $x$ un point de $X$,
$y=f(x)$. Supposons que $\mathcal G$ ait la propriété $(S_k)$ au point $y$, et
que $\mathcal F_y$ ait la propriété $(S_k)$ au point $x$ et soit
équidimensionnel au point $x$. Alors~:
\begin{enumerate}[label=(\roman*)]
  \item $\mathcal F\otimes_{\mathcal O_Y}\mathcal G$ possède la propriété
  $(S_k)$ au point $x$.
  \item Si de plus $Y$ est localement immersible dans un schéma régulier, ou
  est un préschéma excellent (7.8.5), il existe un voisinage de $x$ dans $X$
  tel que $\mathcal F\otimes_{\mathcal O_Y}\mathcal G$ ait la propriété
  $(S_k)$ dans ce voisinage.
\end{enumerate}
\end{corollary}

\begin{proof}
En effet, par (12.1.1, (iv)), il existe un voisinage ouvert $V$ de $x$ tel que
pour tout $x'\in V$, $\mathcal F_{f(x')}$ vérifie $(S_k)$ au point $x'$.
D'autre part, $\mathcal G$ vérifie $(S_k)$ en tous les points de
$Y'=\operatorname{Spec}(\mathcal O_y)$ (5.7.2). Remplaçant $Y$ par $Y'$ et $X$
par $V\times_YY'$, on est ramené (compte tenu de ($\mathbf I$, 3.6.5)) au cas
où $\mathcal G$ vérifie $(S_k)$ dans $Y$ tout entier et où, pour tout
$y\in f(X)$, $\mathcal F_y$ possède la propriété $(S_k)$ en tous les points de
$f^{-1}(y)$~; comme $\mathcal F$ est $f$-plat, on sait alors (6.4.1) que
$\mathcal F\otimes_{\mathcal O_Y}\mathcal G$ possède la propriété $(S_k)$ dans
$X$, ce qui prouve (i). Pour démontrer (ii), il suffit d'observer que, sous les
hypothèses faites, $\mathcal G$ possède la
\oldpage[IV-3]{178}
propriété $(S_k)$ dans un voisinage $U$ de $y$ dans $Y$ ((6.11.2) et (7.8.3,
(iv)))~; on conclut alors de la même manière en remplaçant cette fois $Y$ par
$U$ et $X$ par $V\times_YU$.
\end{proof}

\begin{theorem}[12.1.6]
\label{IV.12.1.6-fr}
Soient $f:X\to Y$ un morphisme plat et localement de présentation finie, $k$
un entier. Les parties suivantes de $X$ sont ouvertes~:
\begin{enumerate}[label=(\roman*)]
  \item L'ensemble des points $x\in X$ tels que $X_{f(x)}$ vérifie la propriété
  $(S_k)$ au point $x$.
  \item L'ensemble des $x\in X$ tels que $X_{f(x)}$ vérifie la propriété
  $(R_k)$ géométrique au point $x$, soit équidimensionnel au point $x$ et tels
  en outre que $x$ n'appartienne à aucun cycle premier immergé associé à
  $X_{f(x)}$.
  \item L'ensemble des $x\in X$ tels que $X_{f(x)}$ soit géométriquement
  régulier (i.e. lisse) au point $x$.
  \item L'ensemble des $x\in X$ tels que $X_{f(x)}$ soit géométriquement
  normal au point $x$.
\end{enumerate}
\end{theorem}

\begin{proof}
Les étapes A) et B) de (12.0.2) se font ici comme dans (12.1.1)~; pour l'étape
A), on utilise (6.7.8), ainsi que (4.2.7) pour (ii)~; pour l'étape B), on
utilise (9.9.2) et (9.9.4). Reste à examiner l'étape C) dans chaque cas.
\begin{enumerate}[label=(\roman*)]
  \item[(i)] Comme dans (12.1.1.2), on se ramène (en utilisant (6.7.8)) au cas
  où $Y$ est le spectre d'un anneau de valuation discrète $A$,
  $X=\operatorname{Spec}(B)$, où $B$ est une $A$-algèbre de type fini, $x$,
  $x'$ deux points de $X$ tels que $f(x)=y$ soit le point fermé et
  $y'=f(x')$ le point générique de $Y$, $x'$ étant en outre une générisation
  de $x$. Comme il s'agit de prouver que $X$ vérifie $(S_k)$ au point $x'$,
  on peut remplacer $X$ par $\operatorname{Spec}(\mathcal O_{X,x})$,
  c'est-à-dire supposer l'anneau $B$ local, l'homomorphisme $A\to B$ étant
  local et $B$ étant une $A$-algèbre essentiellement de type fini (1.3.8).
  Alors, par hypothèse, si $t$ est une uniformisante de $A$, $t$ est un
  élément $B$-régulier par platitude, et $B/tB$ vérifie $(S_k)$. Mais comme
  $A$ est un anneau universellement caténaire (5.6.4), $B$ est caténaire,
  donc, par (5.12.4), il en résulte que $B$ vérifie $(S_k)$, ce qui achève la
  démonstration dans le cas (i).
  \item[(iii)] Ici l'étape C) est inutile~; comme $f$ est plat, on sait en
  effet (6.8.7) que lorsque $Y$ est localement noethérien et $f$ localement de
  type fini, l'ensemble des $x\in X$ tels que $X_{f(x)}$ soit géométriquement
  régulier au point $x$ est ouvert dans $X$.
  \item[(ii)] En raisonnant comme dans (12.1.1.3), pour prouver que la propriété
  considérée est stable par générisation, on peut d'abord, en considérant une
  extension finie quelconque $K'$ de $k(y')$, et tenant compte de la définition
  (6.7.6) de la propriété $(R_k)$ géométrique, ainsi que de l'invariance par
  changement du corps de base des deux autres propriétés figurant dans (ii),
  remplacer dans (ii) la propriété $(R_k)$ géométrique par la propriété
  $(R_k)$. Procédant alors comme dans (i), on se ramène au cas où $Y$ est
  spectre d'un anneau de valuation discrète $A$, et puisque $A$ est caténaire,
  il suffit d'appliquer (5.12.5) pour conclure.
  \item[(iv)] L'ensemble des points $x\in X$ tels que $X_{f(x)}$ soit
  géométriquement normal au point $x$ est contenu dans celui des $x\in X$ tels
  que $X_{f(x)}$ soit géométriquement ponctuellement intègre et vérifie $(S_2)$
  au point $x$, et ce dernier est ouvert dans $X$ en vertu de (12.1.1,
  (viii)). On peut donc déjà supposer que, pour tout $x\in X$, $X_{f(x)}$ est
  géométriquement ponctuellement intègre et vérifie $(S_2)$ au point $x$, et
  \emph{a fortiori} il est équidimensionnel. D'autre part, en vertu du critère
  de Serre (5.8.6), dire que $X_{f(x)}$ est géométriquement normal au point $x$
  signifie que $X_{f(x)}$ vérifie $(S_2)$ et la propriété $(R_1)$ géométrique
  en $x$~;
  \oldpage[IV-3]{179}
  mais en vertu de (ii) et des remarques précédentes, cet ensemble est
  intersection de deux ouverts de $X$, donc est ouvert dans $X$. C.Q.F.D.
\end{enumerate}
\end{proof}

\begin{corollary}[12.1.7]
\label{IV.12.1.7-fr}
Soit $f:X\to Y$ un morphisme localement de présentation finie. Alors
l'ensemble des points $x\in X$ où $f$ possède l'une des propriétés suivantes
(6.8.1)~:
\begin{enumerate}[label=(\roman*)]
  \item vérifier la propriété $(S_k)$~;
  \item être de coprofondeur $\leq n$~;
  \item être de Cohen-Macaulay~;
  \item être régulier (ou lisse, ce qui revient au même)~;
  \item être normal~;
  \item être réduit~;
\end{enumerate}
est ouvert dans $X$.
\end{corollary}

\begin{proof}
En effet, il résulte de (11.3.1) que l'ensemble des $x\in X$ où $f$ est plat
est ouvert. On peut donc se borner au cas où $f$ est plat, et alors le
corollaire résulte de (12.1.1, (iv), (vi) et (vii)) et de (12.1.6, (i), (ii)
et (iv)).
\end{proof}

\begin{remarks}[12.1.8]
\label{IV.12.1.8-fr}
\begin{enumerate}[label=(\roman*)]
  \item Dans (12.1.6, (ii)), on ne peut supprimer l'hypothèse que $x$
  n'appartienne à aucun cycle premier associé immergé de $X_{f(x)}$. C'est ce
  que montre l'exemple (5.12.3), où l'on prend $X=\operatorname{Spec}(A)$, $Y$
  étant le spectre de l'anneau local $A_0$ de $k[T]$ correspondant à l'idéal
  premier $(T)$ et le morphisme $X\to Y$ correspondant à l'homomorphisme
  injectif $A_0\to A$, déduit par localisation et passage au quotient de
  l'injection $k[T]\to k[T,U,V]$~; ce morphisme est plat puisque $A$ est un
  $A_0$-module sans torsion ($\mathbf 0_{\mathrm I}$, 6.3.4). Ici la fibre
  $X_y$ au point fermé $y$ de $Y$, égale à $\operatorname{Spec}(A/tA)$, est
  irréductible, de dimension 1 et vérifie la condition $(R_0)$ géométrique
  puisque l'anneau local en son point générique est un corps. Par contre, au
  point générique $y'$ de $Y$, la fibre $X_{y'}$ a deux composantes
  irréductibles de dimensions respectives 0 et 1, et celle de dimension 0
  n'est pas réduite, donc $X_{y'}$ ne vérifie pas la condition $(R_0)$.
  \item Dans (12.1.6, (ii)), on ne peut pas non plus supprimer l'hypothèse que
  $X_{f(x)}$ est équidimensionnel au point $x$. On le voit ici sur l'exemple
  (5.12.6) avec $X=\operatorname{Spec}(A)$, $Y=\operatorname{Spec}(A_0)$, où
  $A_0$ est défini comme dans (i), le morphisme $X\to Y$ provenant encore de
  l'injection $k[T]\to k[T,U,V,W]$ par localisation et passage au quotient, et
  étant plat puisque $A$ est un $A_0$-module sans torsion
  ($\mathbf 0_{\mathrm I}$, 6.3.4). La fibre $X_y$ au point fermé $y$ de $Y$
  est réduite (donc vérifie $(S_1)$) et vérifie $(R_1)$, mais a deux composantes
  irréductibles de dimensions 2 et 1, tandis que la fibre $X_{y'}$ au point
  générique $y'$ de $Y$ ne vérifie pas $(R_1)$.
\end{enumerate}
\end{remarks}
\subsection{Propriétés locales et globales des fibres d'un morphisme propre,
plat et de présentation finie}
\label{subsection:IV.12.2-fr}

\begin{theorem}[12.2.1]
\label{IV.12.2.1-fr}
Soient $f:X\to Y$ un morphisme \emph{propre} de présentation finie,
$\mathcal F$ un $\mathcal O_X$-Module $f$-plat et de présentation finie,
$\Phi$ une partie finie de $\mathbf Z\cup\{\pm\infty\}$, $k$ un entier. Les
parties suivantes de $Y$ sont ouvertes~:
\begin{enumerate}[label=(\roman*)]
  \item L'ensemble des $y\in Y$ tels que l'ensemble des dimensions des cycles
  premiers associés à $\mathcal F_y$ soit contenu dans $\Phi$.
  \item L'ensemble des $y\in Y$ tels que l'ensemble des dimensions des
  composantes irréductibles de $\operatorname{Supp}(\mathcal F_y)$ contienne
  $\Phi$.
  \item L'ensemble des $y\in Y$ tels que tous les cycles premiers associés à
  $\mathcal F_y$ aient une même dimension égale à $k$.
  \item L'ensemble des $y\in Y$ tels que $\mathcal F_y$ n'ait pas de cycle
  premier associé immergé et que l'ensemble des dimensions des composantes
  irréductibles de $\operatorname{Supp}(\mathcal F_y)$ soit égal à $\Phi$.
  \oldpage[IV-3]{180}
  \item L'ensemble des $y\in Y$ tels que $\mathcal F_y$ ait la propriété
  $(S_k)$ et soit équidimensionnel en tout point de $X_y$.
  \item L'ensemble des $y\in Y$ tels que
  $\operatorname{coprof}(\mathcal F_y)\leq k$.
  \item L'ensemble des $y\in Y$ tels que $\mathcal F_y$ soit un
  $\mathcal O_{X_y}$-Module de Cohen-Macaulay.
  \item L'ensemble des $y\in Y$ tels que $\mathcal F_y$ soit géométriquement
  réduit.
  \item L'ensemble des $y\in Y$ tels que $\mathcal F_y$ soit géométriquement
  ponctuellement intègre en chaque point de $X_y$.
  \item L'ensemble des $y\in Y$ tels que $\mathcal F_y$ soit géométriquement
  intègre.
  \item L'ensemble des $y\in Y$ tels que $\mathcal F_y$ n'ait pas de cycle
  premier associé immergé et que la somme $l(y)$ des multiplicités totales
  (4.7.12) de $\mathcal F_y$ aux points maximaux de
  $\operatorname{Supp}(\mathcal F_y)$ soit $\leq k$.
\end{enumerate}
\end{theorem}

\begin{proof}
À l'exception de (ii), (iii), (iv), (x) et (xi), les propriétés $P(y)$
considérées sont de la forme suivante~: «~pour tout $x\in X_y$, $\mathcal F_y$
a la propriété $Q(x)$~», où $Q(x)$ est l'une des propriétés (i) à (viii) de
(12.1.1). Il résulte de (12.1.1) que l'ensemble $U$ des $x\in X$ tels que
$Q(x)$ soit vraie est ouvert, et l'ensemble $V$ des $y\in Y$ tels que $P(y)$
soit vraie n'est autre que l'ensemble $Y-f(X-U)$~; dans tous ces cas, le
théorème est donc déjà vrai en supposant seulement le morphisme $f$ \emph{fermé}.
Pour (iii), on applique (i) avec $\Phi$ réduit à un seul élément. Il reste à
examiner les cas (ii), (iv) et (xi) séparément, ((x) se déduisant aussitôt de
(xi) et de (4.7.14)) en appliquant toujours la méthode décrite dans (12.0.2).
\end{proof}

\begin{env}[12.2.1.1]
\label{IV.12.2.1.1-fr}
Cas (ii)~: Les étapes A) et B) de la démonstration procèdent comme dans le
début de (12.1.1)~; pour l'étape A), on utilise (8.9.1), (11.2.6) ainsi que
(4.2.7)~; pour l'étape B), on utilise (9.5.5) appliqué à
$\operatorname{Supp}(\mathcal F)$. Il reste à démontrer que si (en supposant
$Y$ noethérien) l'ensemble des dimensions des composantes irréductibles de
$\operatorname{Supp}(\mathcal F_y)$ contient $\Phi$, alors il en est de même
de l'ensemble des dimensions des composantes irréductibles de
$\operatorname{Supp}(\mathcal F_{y'})$, pour toute \emph{générisation} $y'$ de
$y$ dans $Y$. Soit $Y_1$ un spectre d'anneau de valuation discrète tel que, si
$y_1$ et $y'_1$ sont le point fermé et le point générique de $Y_1$, il y ait un
morphisme $h:Y_1\to Y$ tel que $h(y_1)=y$, $h(y'_1)=y'$ (\textbf{II}, 7.1.7).
Appliquant (4.2.7), on voit qu'on peut remplacer $Y$ par $Y_1$ et $X$ par
$X_1=X\times_YY_1$, autrement dit, se borner au cas où $Y$ est spectre
d'anneau de valuation discrète, $y$ son point fermé et $y'$ son point générique.

Utilisant (\textbf{Err}\textsubscript{\textbf{III}}, 30), on peut se borner au
cas où $\operatorname{Supp}(\mathcal F)=X$, de sorte que $f$ est quasi-plat
(2.3.3)~; les composantes irréductibles $Z_i$ de $X$ dominent alors $Y$
(2.3.4), autrement dit leurs points génériques appartiennent à $X_{y'}$, et
les $Z_i\cap X_{y'}$ sont les composantes irréductibles de $X_{y'}$. Mais toute
composante irréductible de $X_y$ est contenue dans un des $Z_i$, donc est une
composante irréductible de $(Z_i)_y$~; or, on sait (7.1.13) que les dimensions
de toutes les composantes irréductibles de $(Z_i)_y$ sont égales à
$\dim((Z_i)_{y'})$, ce qui achève de prouver (ii).
\end{env}

\begin{env}[12.2.1.2]
\label{IV.12.2.1.2-fr}
Cas (iv)~: Le même raisonnement qu'au début montre, en utilisant (12.1.1,
(iii)) que l'on peut déjà supposer (en remplaçant $Y$ par un ouvert de $Y$) que
pour tout $y\in Y$, $\mathcal F_y$ n'a pas de cycle premier associé immergé.
L'assertion du cas (iv) est alors une conséquence immédiate des assertions des
cas (i) et (ii).
\end{env}

\begin{env}[12.2.1.3]
\label{IV.12.2.1.3-fr}
\oldpage[IV-3]{181}
Cas (xi)~: Pour l'étape A) du raisonnement, on utilise (8.9.1), (11.2.6),
(8.10.5, (xii)) (pour conserver l'hypothèse que $f$ est propre), ainsi que
(4.2.7), (4.5.6) et (4.7.9). Pour l'étape B), on voit, comme dans le cas (iv),
qu'on peut supposer que pour tout $y\in Y$, $\mathcal F_y$ n'a pas de cycle
premier associé immergé, et l'on applique (9.8.8) qui prouve que la fonction
$z\mapsto l(z)$ est constructible. Il reste donc à voir (en supposant $Y$
noethérien et $f$ propre) que pour toute générisation $y'$ d'un point $y\in Y$,
on a $l(y')\leq l(y)$. En raisonnant comme dans (12.1.1.3), on voit (à l'aide
de (4.7.8) et (4.5.11)) qu'on peut supposer que les composantes irréductibles
de $\operatorname{Supp}(\mathcal F_{y'})$ sont géométriquement irréductibles
et que $l(y')$ est la \emph{somme des longueurs} des
$(\mathcal F_{y'})_{z'_j}$ aux points maximaux $z'_j$ de
$\operatorname{Supp}(\mathcal F_{y'})$. Appliquant de nouveau (4.2.7), (4.5.6)
et (4.7.9), on se ramène, comme dans (12.2.1.1), au cas où $Y$ est un spectre
d'anneau de valuation discrète, $y$ son point fermé et $y'$ son point générique.
Le fait que $l(y')\leq l(y)$ sera alors conséquence du

\begin{lemma}[12.2.1.4]
\label{IV.12.2.1.4-fr}
Soient $Y$ un spectre d'anneau de valuation discrète, $y$ son point fermé,
$y'$ son point générique, $f:X\to Y$ un morphisme propre, $\mathcal F$ un
$\mathcal O_X$-Module cohérent et $f$-plat, $z_i$ (resp. $z'_j$) les points
maximaux de $\operatorname{Supp}(\mathcal F_y)$ (resp.
$\operatorname{Supp}(\mathcal F_{y'})$). Supposons que $\mathcal F_y$ n'ait
pas de cycle premier associé immergé. Alors on a
\[
  \sum_j\operatorname{long}((\mathcal F_{y'})_{z'_j})
  \leq\sum_i\operatorname{long}((\mathcal F_y)_{z_i}).
\]
\end{lemma}

On a $Y=\operatorname{Spec}(A)$, où $A$ est un anneau de valuation discrète,
dont nous désignerons par $t$ une uniformisante, de sorte que
$\mathcal F_y=\mathcal F/t\mathcal F$. Comme $\mathcal F$ est $f$-plat, les
$z'_j$ sont aussi les points maximaux de $\operatorname{Supp}(\mathcal F)$
(2.3.4)~; pour tout $z_i$, désignons par $z'_{ik}$ ceux des $z'_j$ qui sont des
générisations de $z_i$~; il résulte de (3.4.1.1) que l'on a
\[
  \operatorname{long}((\mathcal F_y)_{z_i})
  \geq\sum_k\operatorname{long}((\mathcal F_{y'})_{z'_{ik}})
\]
d'où en sommant
\[
  \sum_i\operatorname{long}((\mathcal F_y)_{z_i})
  \geq\sum_i\sum_k\operatorname{long}((\mathcal F_{y'})_{z'_{ik}}).
\]

Le lemme sera donc prouvé si nous établissons que pour tout $z'_j$, il y a au
moins un indice $i$ tel que $z'_j$ soit l'un des $z'_{ik}$. Or, comme $f$ est
propre (donc fermé) et $f(z'_j)=y'$, il existe $x\in X$ qui est une
spécialisation de $z'_j$ et est tel que $f(x)=y$, autrement dit
$\overline{\{z'_j\}}\cap X_y$ n'est pas vide. Comme $t$ est
$\mathcal F$-régulier par platitude, on déduit de (3.4.3) qu'il y a au moins un
point de $\operatorname{Ass}(\mathcal F_y)$ dont $z'_j$ est une
générisation~; mais comme les points de $\operatorname{Ass}(\mathcal F_y)$
sont par hypothèse les $z_i$, cela termine la démonstration.
\end{env}

\begin{corollary}[12.2.2]
\label{IV.12.2.2-fr}
Sous les hypothèses de (12.2.1)~:
\begin{enumerate}[label=(\roman*)]
  \item La fonction $y\mapsto\dim(\operatorname{Supp}(\mathcal F_y))$ est
  continue (donc localement constante) dans $Y$.
  \item La fonction $y\mapsto\operatorname{coprof}(\mathcal F_y)$ (5.7.1) est
  semi-continue supérieurement dans $Y$.
  \item La fonction $y\mapsto l(y)$ (12.2.1, (xi)) est semi-continue
  supérieurement dans $Y$ lorsque les $\mathcal F_y$ n'ont pas de cycle
  premier associé immergé.
  \oldpage[IV-3]{182}
  \item La fonction $y\mapsto\operatorname{prof}^{*}(\mathcal F_y)$ (10.8.1)
  est semi-continue inférieurement dans $Y$.
\end{enumerate}
\end{corollary}

\begin{proof}
Il résulte de (12.2.1, (i)) appliqué à $\Phi$ égal au plus petit intervalle de
$\mathbf N$ contenant les dimensions des cycles premiers associés à
$\mathcal F_y$, que $z\mapsto\dim(\operatorname{Supp}(\mathcal F_z))$ est
semi-continue supérieurement au point $y$~; il résulte d'autre part de
(12.2.1, (ii)) appliqué à $\Phi$ égal à l'ensemble des dimensions des
composantes irréductibles de $\operatorname{Supp}(\mathcal F_y)$, que cette
même fonction est semi-continue inférieurement au point $y$, d'où la
conclusion. Les assertions (ii) et (iii) ne sont que d'autres formulations de
(12.2.1, (vi) et (xi)). Enfin, d'après (12.1.3), l'ensemble des $x\in X$ tels
que $\operatorname{prof}^{*}_{f(x)}(\mathcal F_{f(x)})\leq k$ est ouvert, et la
conclusion de (iv) en résulte par le même raisonnement qu'au début de (12.2.1).
\end{proof}

\begin{remarks}[12.2.3]
\label{IV.12.2.3-fr}
\begin{enumerate}[label=(\roman*)]
  \item On observera que pour (12.2.1, (ii)), on peut se dispenser de
  l'hypothèse que $f$ est \emph{propre}.

  \item Dans (12.2.1, (ii)), on ne peut remplacer la condition que l'ensemble
  $E(y)$ des dimensions des composantes irréductibles de
  $\operatorname{Supp}(\mathcal F_y)$ contienne $\Phi$, par la condition que
  $E(y)$ soit \emph{égal à} $\Phi$. En effet, soit $k$ un corps, et considérons
  l'espace projectif à 3 dimensions $P=\mathbf P^3_k=\operatorname{Proj}(C)$,
  où $C=k[t_0,t_1,t_2,t_3]$ ($t_i$ indéterminées)~; dans $P$, soit $X_0$ le
  sous-préschéma fermé «~réunion de la droite $X_1$ définie par $t_1=t_3=0$ et
  du plan $X_2$ défini par $t_2-t_3=0$~», ce qui décrit en termes géométriques
  le fait que $X_0$ est égal à $\operatorname{Proj}(C/\mathfrak a)$, où l'idéal
  gradué $\mathfrak a$ est égal à $\mathfrak b\mathfrak c$, avec
  $\mathfrak b=Ct_1+Ct_3$, $\mathfrak c=C(t_2-t_3)$~; considérons le morphisme
  $f_0:X_0\to X_1$ qui, en termes géométriques, est la «~projection de $X_0$
  sur $X_1$ à partir de la droite à l'infini $D$ définie par $t_0=t_2=0$~»~;
  sous forme algébrique, $f_0$ correspond à l'homomorphisme d'algèbres graduées
  \[
    k[t_0,t_2]\longrightarrow k[t_0,t_1,t_2,t_3]/\mathfrak a
  \]
  obtenu par passage au quotient à partir de l'injection canonique
  $k[t_0,t_2]\to k[t_0,t_1,t_2,t_3]$. Il est clair que $f_0$ est un morphisme
  projectif, donc \emph{propre}~; il en est donc de même de sa restriction
  $f:X\to Y$, où $Y=D_+(t_0)$ et $X=f_0^{-1}(Y)$~; pour voir que $f$ est
  \emph{plat}, il suffit de remarquer que $k[t_2]$ est un anneau principal et
  l'anneau
  \[
    C'=k[t_1,t_2,t_3]/(t_2-t_3)((t_1)+(t_3))
  \]
  un $k[t_2]$-module sans torsion ($\mathbf 0_{\mathrm I}$, 6.3.4). Pour tout
  $y\in Y$ distinct du point $y_0$ défini par $t_2=0$, $f^{-1}(y)$ a alors
  deux composantes irréductibles, de dimensions 0 et 1, tandis que
  $f^{-1}(y_0)$ n'a qu'une seule composante irréductible de dimension 1.

  \item Dans (12.2.1, (i)), on ne peut pas non plus remplacer la condition que
  l'ensemble $E'(y)\supset E(y)$ des dimensions des cycles premiers associés à
  $\mathcal F_y$ soit contenu dans $\Phi$ par la condition qu'il soit
  \emph{égal à} $\Phi$. Soient en effet $k$ un corps, $A$ l'anneau $k[t]$ de
  polynômes, $B$ le sous-anneau $k[t^4,t^3u,tu^3,u^4]$ de l'anneau de polynômes
  $k[t,u]$ ($t,u$ indéterminées)~; $B$ n'est pas intégralement clos, l'élément
  $t^2u^2$ appartenant à sa clôture intégrale mais n'étant pas dans $B$~; si
  $\mathfrak m$ est l'idéal maximal de $B$, engendré par $t^4,t^3u,tu^3$ et
  $u^4$, $\mathfrak m$ est idéal premier associé immergé de $A/At^4$. Posons
  $Y=\operatorname{Spec}(A)$, $X=\operatorname{Spec}(B)$, et soit $f:X\to Y$
  le morphisme correspondant à l'homomorphisme $A\to B$ qui transforme $t$ en
  $t^4$. Comme cet homomorphisme fait de $B$ un $A$-module sans torsion, $f$
  est plat ($\mathbf 0_{\mathrm I}$, 6.3.4). Si $y$ est le point de $Y$
  correspondant à l'idéal maximal $tA$, le préschéma $f^{-1}(y)$ est
  irréductible et de dimension 1, mais admet un cycle premier associé immergé
  de dimension 0~; au contraire, si $y'$ est le point générique de $Y$,
  $f^{-1}(y')$ n'a pas de cycle premier associé immergé, étant le spectre d'une
  $k(y')$-algèbre intègre de dimension 1. Dans cet exemple, $f$ est un morphisme
  affine non propre~; on peut en déduire aussitôt un exemple analogue où $f$
  est propre et plat en considérant $X$ comme un ouvert partout dense d'un
  schéma projectif sur $Y$ (\textbf{II}, 5.3.4 et 5.3.2), ou en procédant
  directement comme dans la Remarque (ii).

  Les Remarques (ii) et (iii) montrent la nécessité d'inclure dans (12.2.1,
  (iv)) la condition que $\mathcal F_y$ n'ait pas de cycle premier associé
  immergé.

  \item Dans (12.2.1, (xi)), on ne peut supprimer l'hypothèse que $f$ soit
  propre. En effet, soient $k$ un corps, $Y$ la «~droite affine~», spectre de
  l'anneau de polynômes $k[t]$ ($t$ indéterminée), $X$ le schéma \emph{somme}
  de $Y$ et de l'ouvert complémentaire du point fermé $y$ de $Y$ défini par
  $t=0$. Il est clair que le morphisme $f:X\to Y$ qui est égal aux injections
  canoniques dans chacune des composantes de $X$ est plat et que si l'on prend
  $\mathcal F=\mathcal O_X$, $\mathcal F_z$ n'a de cycle premier associé
  immergé pour aucun $z\in Y$. Toutefois, on voit aussitôt que l'on a $l(y)=1$
  et $l(z)=2$ pour tout $z\neq y$.

  \item Dans (12.2.1, (xi)), on ne peut pas non plus supprimer l'hypothèse que
  $\mathcal F_y$ soit sans cycle premier associé immergé. C'est ce que montre
  l'exemple de la Remarque (ii)~: on voit aussitôt en effet que l'on a, avec
  $\mathcal F=\mathcal O_X$, $l(y_0)=1$ et $l(y)=2$ pour $y\neq y_0$ dans $Y$.

  \item Nous ignorons si dans (12.2.1, (xi)) on peut remplacer l'hypothèse que
  $\mathcal F$ est $f$-plat par celle que $\operatorname{Supp}(\mathcal F)=X$
  et que $f$ est universellement ouvert. En reprenant la démonstration de
  (12.2.1.4), on voit (en utilisant aussi (14.3.6)) qu'il faudrait résoudre la
  question suivante~: Soient $A$ un anneau local noethérien, $M$ un $A$-module
  \oldpage[IV-3]{183}
  de type fini, $t$ un élément de l'idéal maximal $\mathfrak m$ de $A$, qui
  n'est contenu dans aucun idéal premier minimal de $A$~; s'il existe un de ces
  idéaux premiers minimaux $\mathfrak p$ tels que $\dim(A/\mathfrak p)=1$,
  est-il vrai que $\mathfrak m$ est un idéal premier associé à $M/tM$ ($t$
  n'étant pas supposé $M$-régulier)~?

  On notera toutefois qu'on ne peut, dans (12.2.1, (xi)), supprimer purement et
  simplement l'hypothèse que $\mathcal F$ est $f$-plat. On prendra en effet ici
  $P$ comme dans la Remarque (ii), et dans $P$ le sous-préschéma fermé $X_0$
  réunion des trois droites $X_1,X_2,X_3$ définies respectivement par
  $t_1=t_3=0$, $t_1-t_0=t_3=0$, $t_2=t_3=0$~; on définit la projection $f_0$ de
  $X_0$ sur $X_1$ comme précédemment et sa restriction $f:X\to Y$, où
  $Y=D_+(t_0)$ est la droite affine et $X=f_0^{-1}(Y)$~; $f$ est propre mais
  non plat, et si l'on prend $\mathcal F=\mathcal O_X$, $\mathcal F_y$ n'a de
  cycles premiers associés immergés pour aucun $y\in Y$. Mais on a $l(y_0)=1$
  et $l(y)=2$ pour $y\neq y_0$ dans $Y$.
\end{enumerate}
\end{remarks}

\begin{theorem}[12.2.4]
\label{IV.12.2.4-fr}
Soient $f:X\to Y$ un morphisme propre, plat et de présentation finie, $k$ un
entier $\geq 1$. Les parties suivantes de $Y$ sont ouvertes~:
\begin{enumerate}[label=(\roman*)]
  \item L'ensemble des $y\in Y$ tels que $X_y$ possède la propriété $(S_k)$.
  \item L'ensemble des $y\in Y$ tels que $X_y$ vérifie la propriété $(R_k)$
  géométrique, soit équidimensionnel en chaque point et n'ait pas de cycle
  premier associé immergé.
  \item L'ensemble des $y\in Y$ tels que $X_y$ soit géométriquement régulier
  (i.e. lisse sur $k(y)$).
  \item L'ensemble des $y\in Y$ tels que $X_y$ soit géométriquement normal.
  \item L'ensemble des $y\in Y$ tels que $X_y$ soit géométriquement réduit.
  \item L'ensemble des $y\in Y$ tels que $X_y$ soit géométriquement réduit et
  que le nombre géométrique de composantes connexes de $X_y$ soit égal à $k$.
  \item L'ensemble des $y\in Y$ tels que $X_y$ soit géométriquement
  ponctuellement intègre (4.6.9).
  \item L'ensemble des $y\in Y$ tels que $X_y$ soit géométriquement intègre.
  \item L'ensemble des $y\in Y$ tels que $X_y$ n'ait pas de cycle premier
  associé immergé et que la multiplicité totale (4.7.4) de $X_y$ sur $k(y)$
  soit $\leq k$.
\end{enumerate}
\end{theorem}

\begin{proof}
Sauf pour (vi), ces assertions sont des cas particuliers d'assertions de
(12.2.1), ou se déduisent d'assertions de (12.1.6) comme au début de la
démonstration de (12.2.1). Pour (vi), on se ramène comme toujours (compte tenu
de l'invariance du nombre géométrique de composantes connexes par changement
de base (4.5.6)) au cas où $Y$ est noethérien. L'ensemble $U$ des $y\in Y$ tels
que $X_y$ soit géométriquement réduit est ouvert dans $Y$ en vertu de (v)~; il
résulte alors de (\textbf{III}, 7.8.7 et 7.8.6) que pour tout $y\in U$, il y a
un voisinage $V\subset U$ de $y$ et un entier $m$ tels que, pour \emph{tout}
$z\in V$, $\Gamma(X_z,\mathcal O_{X_z})$ soit isomorphe à $(k(z))^m$~; mais en
vertu de (\textbf{III}, 4.3.4), $m$ est alors le nombre géométrique de
composantes connexes de $X_z$, d'où la conclusion. On donnera une autre
démonstration de (vi) dans (15.5.9).
\end{proof}

\subsection{Propriétés cohomologiques locales des fibres d'un morphisme plat et
localement de présentation finie}
\label{subsection:IV.12.3-fr}

\begin{lemma}[12.3.1]
\label{IV.12.3.1-fr}
Soient $A$ un anneau, $B$ une $A$-algèbre de présentation finie qui soit un
$A$-module plat, $M$ un $B$-module de présentation finie, qui soit un
$A$-module plat. Alors les conditions suivantes sont équivalentes~:
\begin{enumerate}[label=\emph{\alph*)}]
  \item $M$ est un $B$-module projectif.
  \item $M$ est un $B$-module plat.
  \item Pour tout $y\in\operatorname{Spec}(A)$, $M\otimes_Ak(y)$ est un
  $(B\otimes_Ak(y))$-module projectif.
\end{enumerate}
\end{lemma}

\begin{proof}
\oldpage[IV-3]{184}
L'équivalence de $a)$ et $b)$ résulte de Bourbaki, \emph{Alg. comm.}, chap. II,
\S~5, n\textsuperscript{o} 2, cor. 2 du th. 1. Comme $a)$ implique $c)$
trivialement, il reste à prouver que $c)$ implique $b)$, ce qui résulte du
critère de platitude par fibres (11.3.10), appliqué avec $g=h$, $f=1_X$.
\end{proof}

\begin{proposition}[12.3.2]
\label{IV.12.3.2-fr}
Soient $A$ un anneau, $B$ une $A$-algèbre de présentation finie qui soit un
$A$-module plat, $M$ un $B$-module de présentation finie qui soit un
$A$-module plat. Alors~:
\begin{enumerate}[label=(\roman*)]
  \item Il existe une résolution gauche de $M$ par des $B$-modules libres de
  type fini.
  \item On a
  \begin{equation}
    \tag{12.3.2.1}
    \label{IV.12.3.2.1-fr}
    \operatorname{dim.proj}_B(M)
    =\sup_{y\in\operatorname{Spec}(A)}
    \operatorname{dim.proj}_{B\otimes_Ak(y)}(M\otimes_Ak(y))
    =\operatorname{Tor.dim}_B(M)
  \end{equation}
  où $\operatorname{Tor.dim}_B(M)$ est le \emph{plus petit entier} $i$ tel que
  $\operatorname{Tor}^B_j(M,N)=0$ pour tout $j\geq i$ et tout $B$-module $N$
  (et $+\infty$ si un tel entier $i$ n'existe pas).
\end{enumerate}
\end{proposition}

\begin{proof}
\begin{enumerate}[label=(\roman*)]
  \item En vertu de (8.9.1), il existe un sous-anneau noethérien $A_0$ de $A$,
  une $A_0$-algèbre de type fini $B_0$ et un $B_0$-module de type fini $M_0$
  tels que $B$ soit isomorphe à $B_0\otimes_{A_0}A$ et $M$ à
  $M_0\otimes_{A_0}A$. De plus (11.2.7) on peut supposer que $B_0$ et $M_0$
  sont des $A_0$-modules plats. Il y a alors une résolution gauche
  $(L_0)_\bullet$ de $M_0$ formée de $B_0$-modules libres de type fini, et
  comme $M_0$ et les $L_{0i}$ sont des $A_0$-modules plats,
  $(L_0)_\bullet\otimes_{A_0}A$ est une résolution gauche de $M$ formée de
  $B$-modules libres de type fini (2.1.10).

  \item En vertu de ($\mathbf 0$, 17.2.2, (ii)), on a
  $\operatorname{Tor.dim}_B(M)\leq\operatorname{dim.proj}_B(M)$, et la
  définition de la dimension projective prouve aussitôt que, pour tout
  $y\in\operatorname{Spec}(A)$, on a
  \[
    \operatorname{dim.proj}_{B\otimes_Ak(y)}(M\otimes_Ak(y))
    \leq\operatorname{dim.proj}_B(M).
  \]
  Pour prouver les inégalités inverses, considérons une résolution gauche
  $(L_i)$ de $M$ par des $B$-modules libres de type fini, et supposons que
  $\operatorname{Tor.dim}_B(M)=n$ (resp.
  $\operatorname{dim.proj}_{B\otimes_Ak(y)}(M\otimes_Ak(y))\leq n$ pour tout
  $y\in\operatorname{Spec}(A)$). Alors
  $R=\operatorname{Im}(L_n\to L_{n-1})=\operatorname{Ker}(L_{n-1}\to L_{n-2})$
  est un $B$-module de type fini et aussi un $A$-module plat, en vertu de
  l'hypothèse sur $M$ et $B$ et de (2.1.10). En outre, on a
  $\operatorname{Tor}^B_{n+1}(M,N)=\operatorname{Tor}^B_1(R,N)$ pour tout
  $B$-module $N$ (M, V, 7). L'hypothèse $\operatorname{Tor.dim}_B(M)=n$
  entraîne donc $\operatorname{Tor}^B_1(R,N)=0$ pour tout $B$-module $N$,
  c'est-à-dire que $R$ est un $B$-module plat, donc projectif en vertu de
  (12.3.1)~; cela établit que $\operatorname{dim.proj}_B(M)\leq n$.
  L'hypothèse
  $\operatorname{dim.proj}_{B\otimes_Ak(y)}(M\otimes_Ak(y))\leq n$ pour tout
  $y\in\operatorname{Spec}(A)$ entraîne d'autre part, par tensorisation avec
  $k(y)$, que dans chacune des suites (exactes en vertu de la platitude sur $A$
  de $M$, des $L_i$ et de $R$ (2.1.10))
  \[
    0\to R\otimes_Ak(y)\to L_{n-1}\otimes_Ak(y)\to\cdots
    \to L_0\otimes_Ak(y)\to M\otimes_Ak(y)\to0
  \]
  $R\otimes_Ak(y)$ est un $(B\otimes_Ak(y))$-module projectif (pour
  $y\in\operatorname{Spec}(A)$). On conclut alors encore de (12.3.1) que $R$
  est un $B$-module projectif, donc $\operatorname{dim.proj}_B(M)\leq n$.
\end{enumerate}
\end{proof}

\begin{proposition}[12.3.3]
\label{IV.12.3.3-fr}
Soient $f:X\to S$ un morphisme localement de présentation finie,
$\mathcal L_\bullet$ un complexe formé de $\mathcal O_X$-Modules
quasi-cohérents de présentation finie~; pour tout $s\in S$, soit
$(\mathcal L_\bullet)_s$ le complexe
$\mathcal L_\bullet\otimes_{\mathcal O_S}k(s)$ de $\mathcal O_{X_s}$-Modules de
type fini. Supposons que $\mathcal L_{n-1}$ soit $f$-plat. Alors l'ensemble $U$
des $x\in X$ tels que $(\mathcal H_n((\mathcal L_\bullet)_{f(x)}))_x=0$ est
ouvert dans $X$. Si de plus $\mathcal L_n$ est $f$-plat, alors on a
$\mathcal H_n(\mathcal L_\bullet)|U=0$.
\end{proposition}

\begin{proof}
On peut évidemment se limiter au cas où $n=1$ et où le complexe
$\mathcal L_\bullet$ se réduit à
$0\to\mathcal L_2\xrightarrow{u}\mathcal L_1\xrightarrow{v}\mathcal L_0\to0$.
On peut d'abord se ramener au cas où $S=\operatorname{Spec}(A)$ et $X$
\oldpage[IV-3]{185}
sont affines, puis au cas où $S$ est \emph{noethérien}~; en effet, par (8.9.1)
et (8.5.2), on sait qu'il existe un préschéma noethérien
$S'=\operatorname{Spec}(A')$, où $A'$ est un sous-anneau de $A$, un morphisme
de type fini $f':X'\to S'$ et trois $\mathcal O_{X'}$-Modules cohérents
$\mathcal L'_i$ ($0\leq i\leq2$) tels que $X=X'\times_{S'}S$,
$f=f'_{(S)}$, $\mathcal L_i=\mathcal L'_i\otimes_{S'}S$, ainsi que deux
homomorphismes $u'$, $v'$ tels que $u=u'\otimes1$, $v=v'\otimes1$ et
$v'\circ u'=0$. On peut en outre supposer que $\mathcal L'_0$ est $f'$-plat
(11.2.7), et que $\mathcal L'_1$ est $f'$-plat lorsque $\mathcal L_1$ est
supposé $f$-plat. Par fidèle platitude, l'hypothèse
$(\mathcal H_n((\mathcal L_\bullet)_{f(x)}))_x=0$ équivaut à
$(\mathcal H_n((\mathcal L'_\bullet)_{f'(x')}))_{x'}=0$, si $x'$ est la
projection de $x$ dans $X'$~; si $U'$ est l'ensemble des $x'\in X'$ tels que
$(\mathcal H_n((\mathcal L'_\bullet)_{f'(x')}))_{x'}=0$, $U$ est donc l'image
réciproque de $U'$ par la projection $X\to X'$, ce qui ramène, pour la première
assertion, au cas où $S$ est noethérien. Pour la seconde assertion, on remarque
en outre que $A$ est limite inductive de ses sous-anneaux $A_\lambda$ qui sont
des $A'$-algèbres de type fini~; posons
$X_\lambda=X'\times_{S'}S_\lambda$, où
$S_\lambda=\operatorname{Spec}(A_\lambda)$,
$\mathcal L_i^{(\lambda)}=\mathcal L'_i\otimes_{S'}S_\lambda$, et soient
$u_\lambda=u'\otimes1:\mathcal L_2^{(\lambda)}\to\mathcal L_1^{(\lambda)}$,
$v_\lambda=v'\otimes1:\mathcal L_1^{(\lambda)}\to\mathcal L_0^{(\lambda)}$,
de sorte que l'on a
$\mathcal H_1(\mathcal L_\bullet^{(\lambda)})
=\operatorname{Ker}(v_\lambda)/\operatorname{Im}(u_\lambda)$. Or, comme le
foncteur $\varinjlim$ est exact dans la catégorie des groupes commutatifs, on a
$\operatorname{Ker}(v)=\varinjlim\operatorname{Ker}(v_\lambda)$,
$\operatorname{Im}(u)=\varinjlim\operatorname{Im}(u_\lambda)$, et
$\mathcal H_1(\mathcal L_\bullet)
=\varinjlim\mathcal H_1(\mathcal L_\bullet^{(\lambda)})$. Si l'on a supposé
que $\mathcal L_1$ est $f$-plat et (en se ramenant au cas où $X=U$) que l'on
ait prouvé $\mathcal H_1(\mathcal L_\bullet^{(\lambda)})=0$ pour tout
$\lambda$, on en déduira bien $\mathcal H_1(\mathcal L_\bullet)=0$.

\medskip
I) Supposons désormais $S$ noethérien. On sait (sans hypothèse de platitude sur
$\mathcal L_0$) que l'ensemble $U$ est \emph{constructible} dans $X$ (9.9.6).
Utilisant maintenant ($\mathbf 0_{\mathrm{III}}$, 9.2.5), il reste à montrer
que pour toute \emph{générisation} $x'$ de $x\in U$ dans $X$, on a aussi
$x'\in U$. La méthode exposée dans (12.0.2) s'applique sans changement (compte
tenu du fait que pour un préschéma $Z$ sur un corps $k$ et une extension $k'$
de $k$, la projection $Z\otimes_k k'\to Z$ est fidèlement plate). On peut donc
supposer que $S$ est le spectre d'un anneau de valuation discrète $A$, dont on
désigne par $t$ une uniformisante, $x$ étant au-dessus du point fermé de $S$ et
$x'$ au-dessus du point générique de $S$. L'hypothèse que $\mathcal L_0$ est
$f$-plat entraîne que $t$ est $\mathcal L_0$-régulier
($\mathbf 0_{\mathrm I}$, 6.3.4). On est alors ramené à prouver le lemme suivant
(où $B$, $M$, $N$, $P$ seront remplacés par $\mathcal O_x$,
$(\mathcal L_2)_x$, $(\mathcal L_1)_x$ et $(\mathcal L_0)_x$ respectivement)~:

\begin{lemma}[12.3.3.1]
\label{IV.12.3.3.1-fr}
Soient $B$ un anneau local noethérien, $M,N,P$ trois $B$-modules, $N$ étant
supposé de type fini, $M\xrightarrow{u}N\xrightarrow{v}P$ deux homomorphismes
tels que $v\circ u=0$. Soit $t$ un élément de l'idéal maximal de $B$ tel que
$t$ soit $P$-régulier et que la suite
\begin{equation}
  \tag{12.3.3.2}
  \label{IV.12.3.3.2-fr}
  M/tM\xrightarrow{u\otimes1_{B/tB}}N/tN
  \xrightarrow{v\otimes1_{B/tB}}P/tP
\end{equation}
soit exacte. Alors la suite $M\xrightarrow{u}N\xrightarrow{v}P$ est exacte.
\end{lemma}

Notons d'abord que si l'on remplace $M$ par son image $Q$ dans $N$ et $u$ par
l'injection $j:Q\to N$, l'image de $j\otimes1$ est la même que celle de
$u\otimes1$, donc on peut, sans changer ni l'hypothèse ni la conclusion,
supposer $u$ injectif. D'autre part, si $R$ est l'image de $v$ et
$\rho:N\to R$ la surjection canonique, $t$ est évidemment $R$-régulier, on a
$\rho\circ u=0$, et le noyau de $\rho\otimes1$ est contenu dans celui de
$v\otimes1$~; il lui est par suite égal si la suite (12.3.3.2) est exacte~;
comme d'autre part on a évidemment $\operatorname{Ker}(\rho)
=\operatorname{Ker}(v)$, on voit qu'on peut, pour prouver le lemme, supposer en
outre $v$ surjectif. Le lemme sera alors conséquence du suivant~:

\begin{lemma}[12.3.3.3]
\label{IV.12.3.3.3-fr}
Soient $B$ un anneau, $M,N,P$ trois $B$-modules,
$u:M\to N$, $v:N\to P$ deux homomorphismes tels que $u$ soit injectif, $v$
surjectif et $v\circ u=0$. Soit $t$ un élément de $B$ tel que $t$ soit
$P$-régulier. Alors on a
\begin{equation}
  \tag{12.3.3.4}
  \label{IV.12.3.3.4-fr}
  \operatorname{Ker}(v\otimes1)/\operatorname{Im}(u\otimes1)
  =(\operatorname{Ker}(v)/\operatorname{Im}(u))\otimes_B(B/tB)
\end{equation}
à un isomorphisme canonique près.
\end{lemma}

En effet, l'hypothèse que la suite (12.3.3.2) est exacte entraînera alors
$(\operatorname{Ker}v/\operatorname{Im}u)\otimes_B(B/tB)=0$, et comme $t$
appartient à l'idéal maximal de $B$ et que $\operatorname{Ker}(v)$ est un
$B$-module de type fini (puisque $B$ est supposé noethérien dans (12.3.3.1)),
le lemme de Nakayama prouvera que $\operatorname{Ker}(v)=\operatorname{Im}(u)$.

Reste donc à démontrer (12.3.3.3). Posons $u'=u\otimes1$, $v'=v\otimes1$,
$Z=\operatorname{Ker}(v)$, $H=Z/\operatorname{Im}(u)$, de sorte qu'on a les
suites exactes
\begin{gather*}
  0\to M\to Z\to H\to0,\\
  0\to Z\to N\xrightarrow{v}P\to0,
\end{gather*}
\oldpage[IV-3]{186}
d'où, en tensorisant par $B/tB$ et utilisant le lemme (3.4.1.4) et le fait que
$t$ est $P$-régulier, les suites exactes
\begin{gather*}
  M/tM\xrightarrow{w'}Z/tZ\to H/tH\to0,\\
  0\to Z/tZ\to N/tN\xrightarrow{v'}P/tP\to0,
\end{gather*}
d'où $\operatorname{Ker}(v')=Z/tZ$, et comme
$\operatorname{Im}(w')=\operatorname{Im}(u')$, on obtient (12.3.3.4).

\medskip
II) Supposons maintenant en outre que $\mathcal L_1$ soit $f$-plat~;
remplaçant en outre éventuellement $X$ par $U$, on peut supposer que $X=U$.
Il s'agit de voir que pour tout $x\in X$, on a
$(\mathcal H_1(\mathcal L_\bullet))_x=0$. Soit $s=f(x)$, et soit
$\mathfrak n$ l'idéal $\mathfrak m_s\mathcal O_{X,x}$ qui est contenu dans
l'idéal maximal de $\mathcal O_{X,x}$~; comme
$(\mathcal H_1(\mathcal L_\bullet))_x$ est un $\mathcal O_{X,x}$-module de type
fini ($X$ étant localement noethérien), il est séparé pour la topologie
$\mathfrak n$-adique ($\mathbf 0_{\mathrm I}$, 7.3.5), donc il suffit de
montrer que son \emph{séparé complété} pour cette topologie est 0. Or, en vertu
de (\textbf{III}, 7.4.7.2), ce séparé complété est égal à
$\varprojlim_k\mathcal H_1(\mathcal L_\bullet\otimes_{\mathcal O_S}
(\mathcal O_{S,s}/\mathfrak m_s^{k+1}))$, et il suffira donc de prouver que
chacun des termes de ce système projectif est nul. Mais cela est vrai par
hypothèse pour $k=0$~; raisonnons donc par récurrence sur $k$. La conclusion
résultera du lemme plus général suivant (qu'on appliquera pour
$A=\mathcal O_{S,s}/\mathfrak m_s^{k+1}$ et $\mathfrak J$ égal à l'idéal
maximal $\mathfrak m_s/\mathfrak m_s^{k+1}$ de $A$)~:

\begin{lemma}[12.3.3.5]
\label{IV.12.3.3.5-fr}
Soient $A$ un anneau, $\mathfrak J$ un idéal nilpotent de $A$,
$L_\bullet:P\xrightarrow{u}Q\xrightarrow{v}R$ un complexe de $A$-modules. On
pose $A_0=A/\mathfrak J$,
$L_\bullet^{(0)}=L_\bullet\otimes_AA_0:
P_0\xrightarrow{u_0}Q_0\xrightarrow{v_0}R_0$, et l'on suppose que
$\operatorname{gr}_{\mathfrak J}^{\bullet}(A)$ est un $A_0$-module plat, et
que $Q$ et $R$ sont des $A$-modules plats. Alors la relation
$\mathrm H_1(L_\bullet^{(0)})=0$ entraîne $\mathrm H_1(L_\bullet)=0$.
\end{lemma}

Posons
$L_\bullet^{(k)}=L_\bullet\otimes_A(A/\mathfrak J^{k+1}):
P_k\xrightarrow{u_k}Q_k\xrightarrow{v_k}R_k$~; comme il existe un entier
$k\geq1$ tel que $L_\bullet^{(k)}=L_\bullet$, il suffira de prouver, par
récurrence sur $k$, que la suite
$P_k\xrightarrow{u_k}Q_k\xrightarrow{v_k}R_k$ est
\oldpage[IV-3]{187}
exacte (cette assertion découlant de l'hypothèse pour $k=0$). Soit donc
$x\in Q_k$ un élément tel que $v_k(x)=0$. Comme par l'hypothèse de récurrence
la suite
$P_{k-1}\xrightarrow{u_{k-1}}Q_{k-1}\xrightarrow{v_{k-1}}R_{k-1}$ est exacte,
l'image canonique de $x$ dans
$Q_{k-1}=Q_k/(\mathfrak J^kQ/\mathfrak J^{k+1}Q)$ appartient à
$\operatorname{Im}(u_{k-1})$, donc il existe $z\in P_k$ tel que
$x=u_k(z)+x'$ avec $x'\in\mathfrak J^kQ/\mathfrak J^{k+1}Q$. Comme on a
$v_k\circ u_k=0$, la relation $v_k(x)=0$ entraîne $v_k(x')=0$, et par ailleurs
on a évidemment $v_k(x')\in\mathfrak J^kR/\mathfrak J^{k+1}R$~; tout revient
donc à prouver que la suite
\[
  \mathfrak J^kP/\mathfrak J^{k+1}P\xrightarrow{u'}
  \mathfrak J^kQ/\mathfrak J^{k+1}Q\xrightarrow{v'}
  \mathfrak J^kR/\mathfrak J^{k+1}R
\]
(où $u'$ et $v'$ proviennent de $u$ et $v$ par restriction et passage aux
quotients) est exacte. Or, par hypothèse,
$\mathfrak J^k/\mathfrak J^{k+1}$ est un $(A/\mathfrak J)$-module plat, donc la
suite
\[
  (P/\mathfrak JP)\otimes_{A/\mathfrak J}
  (\mathfrak J^k/\mathfrak J^{k+1})\xrightarrow{u''}
  (Q/\mathfrak JQ)\otimes_{A/\mathfrak J}
  (\mathfrak J^k/\mathfrak J^{k+1})\xrightarrow{v''}
  (R/\mathfrak JR)\otimes_{A/\mathfrak J}
  (\mathfrak J^k/\mathfrak J^{k+1})
\]
est exacte. Mais
$(Q/\mathfrak JQ)\otimes_{A/\mathfrak J}
(\mathfrak J^k/\mathfrak J^{k+1})
=Q\otimes_A(\mathfrak J^k/\mathfrak J^{k+1})$ s'identifie à
$\mathfrak J^kQ/\mathfrak J^{k+1}Q$ en vertu de la platitude de $Q$ sur $A$,
et de même
$(R/\mathfrak JR)\otimes_{A/\mathfrak J}
(\mathfrak J^k/\mathfrak J^{k+1})$ s'identifie à
$\mathfrak J^kR/\mathfrak J^{k+1}R$. Enfin, l'image de $u'$ s'identifie à
celle de $u''$, et cela termine la démonstration.
\end{proof}

\begin{corollary}[12.3.4]
\label{IV.12.3.4-fr}
Soient $f:X\to S$ un morphisme plat localement de présentation finie,
$\mathcal F$, $\mathcal G$ deux $\mathcal O_X$-Modules de présentation finie et
$f$-plats. Soient $n$ un entier $\geq0$, $U$ l'ensemble des $x\in X$ tels que
\[
  \bigl(\mathscr{E}\!xt^n_{\mathcal O_{X_{f(x)}}}
  (\mathcal F_{f(x)},\mathcal G_{f(x)})\bigr)_x=0
\]
(resp.
$\bigl(\mathscr{T}\!or_n^{\mathcal O_{X_{f(x)}}}
(\mathcal F_{f(x)},\mathcal G_{f(x)})\bigr)_x=0$). Alors $U$ est ouvert, et
l'on a
\[
  \mathscr{E}\!xt^n_{\mathcal O_X}(\mathcal F,\mathcal G)|U=0
\]
(resp.
$\mathscr{T}\!or_n^{\mathcal O_X}(\mathcal F,\mathcal G)|U=0$).
\end{corollary}

\begin{proof}
On peut évidemment se borner au cas où $S$ et $X$ sont affines. Alors (12.3.2),
il existe une résolution gauche $\mathcal L_\bullet=(\mathcal L_i)$ de
$\mathcal F$ formée de $\mathcal O_X$-Modules \emph{libres de type fini}. Si
l'on pose $\mathcal M^\bullet=\mathscr{H}\!om(\mathcal L_\bullet,\mathcal G)$
(resp. $\mathcal N_\bullet=\mathcal L_\bullet\otimes_{\mathcal O_X}\mathcal G$),
on en déduit que chacun des $\mathcal M^i$ (resp. $\mathcal N_i$) est isomorphe
à un $\mathcal O_X$-Module de la forme $\mathcal G^m$, donc les $\mathcal M^i$
et $\mathcal N_i$ sont des $\mathcal O_X$-Modules de présentation finie et
$f$-plats. En outre, pour tout changement de base $S'\to S$, si l'on pose
$X'=X\times_SS'$, $\mathcal F'=\mathcal F\otimes_SS'$,
$\mathcal G'=\mathcal G\otimes_SS'$, $\mathcal L'_\bullet
=\mathcal L_\bullet\otimes_SS'$, $\mathcal L'_\bullet$ est encore une
résolution gauche de $\mathcal F'$ par des $\mathcal O_{X'}$-Modules libres de
type fini (2.1.10), et $\mathcal M'^\bullet=\mathcal M^\bullet\otimes_SS'$ est
égal à $\mathscr{H}\!om(\mathcal L'_\bullet,\mathcal G')$ (resp.
$\mathcal N'_\bullet=\mathcal N_\bullet\otimes_SS'$ est égal à
$\mathcal L'_\bullet\otimes_{\mathcal O_{X'}}\mathcal G'$), d'après ce qui
précède. En particulier, on a, pour tout $s\in S$,
\[
  \mathscr{E}\!xt^n_{\mathcal O_{X_s}}(\mathcal F_s,\mathcal G_s)
  =\mathcal H^n((\mathcal M^\bullet)_s)
\]
(resp.
$\mathscr{T}\!or_n^{\mathcal O_{X_s}}(\mathcal F_s,\mathcal G_s)
=\mathcal H_n((\mathcal N_\bullet)_s)$). Appliquant (12.3.3) aux complexes de
Modules $f$-plats $\mathcal M^\bullet$ et $\mathcal N_\bullet$, on en déduit
aussitôt le corollaire.
\end{proof}

% IV3_B26_P187_END_BEFORE_SECTION_13
